%%%----------------------------------------------------------------------------------------
%	SECTION TITLE
%----------------------------------------------------------------------------------------
\vspace{-3mm}
\cvsection{Experience}

%----------------------------------------------------------------------------------------
%	SECTION CONTENT
%----------------------------------------------------------------------------------------

\begin{cventries}

%------------------------------------------------
\cventry
{Staff Engineer II (Engineering Physicist II)}
{SLAC National Laboratory}
{Palo Alto, CA}
{April 2022 - Present}
{
\begin{cvitems}
\item {Led the restoration and recommissioning of the FACET positron beamline (inactive since 2016), returning RF, magnet, vacuum, and water cooling systems to operational status.}
\item {Directed cross-disciplinary engineering efforts across vacuum, electrical, cooling, and beam diagnostic systems to re-establish beamline functionality and ensure safe, reliable operation.}
\item {Led the recommissioning of the bunch length diagnostic system after six years offline, restoring a critical beam diagnostic and enabling reliable, high-resolution measurements for accelerator operations and research.}
\item {Developed MATLAB-based tools to process and visualize beam diagnostics data across accelerator sectors, including magnet mover controls and slice emittance and energy spread analysis, improving operational diagnostics and analytical capability for accelerator studies.}
\item {Supported FACET-II accelerator operations and troubleshooting efforts, diagnosing issues across RF, magnet, vacuum, and diagnostic systems to maintain stable beam delivery and operational reliability.}
\item {Engineered electrical configurations to enforce robust Lockout-Tagout (LOTO) conditions, significantly improving equipment safety, maintenance accessibility, and operational reliability.}
\end{cvitems}
}
\vspace{6mm}
%------------------------------------------------    
\cventry
{Engineering Physicist I, US-HL-AUP Superconducting R\&D Engineer \& Electrical Quality Assurance}
{Fermilab}
{Batavia, IL}
{April 2015 - March 2022}
{
\begin{cvitems}
\item {Held primary responsibility for electrical quality control of Nb$_3$Sn superconducting magnet systems in the US HiLumi-LHC AUP, improving fabrication reliability and reducing early-stage production defects.}
\item {Performed electromechanical failure analyses on prototype and production Nb$_3$Sn coils to diagnose quench mechanisms and inform time-critical production decisions.}
\item {Designed and fabricated small-scale R\&D superconducting test coils to reproduce quench failure modes observed in pre-production magnet coils, supporting electromechanical failure investigations.}
\item {Designed and implemented a nitrogen (N$_2$) distribution system for furnace operations to control oxygen during Nb$_3$Sn coil processing, improving process reliability and reducing technician intervention.}
\item {Developed a MATLAB GUI to analyze thermocouple temperature histories against heat treatment recipes, identifying process control discrepancies and supporting furnace performance validation.}
\item {Managed daily fabrication workflow for Nb$_3$Sn superconducting magnet coils, supporting pre-production (24 units) and full production (50 units) while coordinating manufacturing priorities to meet delivery schedules.}
\item {Collaborated with physicists and engineers to optimize Nb$_3$Sn coil heat treatment processes and achieve targeted superconducting performance metrics based on laboratory testing and analysis.}
\item {Executed final electrical qualification testing of AUP superconducting coils prior to shipment to CERN, verifying compliance with specifications and readiness for system integration.}
\end{cvitems}
}
\vspace{3mm}
%------------------------------------------------

\cventry
{Accelerator Systems \& Operations}
{Fermilab}
{Batavia, IL}
{Feb 2013 - April 2014}
{
\begin{cvitems}
\item {Coordinated accelerator operations to initiate, optimize, and maintain beamline performance in support of laboratory and university-based research programs.}
\item {Operated, tuned, and maintained the H$^-$ plasma source to ensure stable beam delivery across the accelerator complex.}
\item {Contributed to the design and implementation of an air diversion system for the NuMI horn power stripline, improving thermal management and operational reliability for high-power beam delivery.}
\end{cvitems}
}
\vspace{6mm}
%------------------------------------------------
\cventry
{Graduate Research Assistant}
{Georgia Institute of Technology}
{Atlanta, GA}
{Dec 2013 - April 2014}
{
\begin{cvitems}
\item {Participated in an in-depth study to understand plasma boundary interactions and wall effects.}
\item {Studied plume divergence behavior of the P5 5 kW Hall thruster.}
\item {Assisted in the design of a glass injection tube for a helicon ion thruster.}
\end{cvitems}
}
\vspace{6mm}
%------------------------------------------------
\cventry
{Active Aeroelasticity and Structures Lab - Undergraduate Research Assistant}
{University of Michigan}
{Ann Arbor, MI}
{Dec 2012 - Dec 2013}
{
\begin{cvitems}
\item {Conducted aeroelasticity research on highly flexible aircraft, supporting design and testing of UAV systems.}
\item {Designed and validated a ventral fin using CFD and wind tunnel testing (2'×2' and 5'×7') to improve yaw stability and control authority.}
\item {Developed a sensor-based system using gyroscopes to model and measure aircraft structural deformation in flight.}
\item {Fabricated and maintained structural components for the X-HALE UAV, supporting iterative design improvements.}
\item {Contributed to full-scale assembly of an Aerial Test Vehicle (ATV), including composite wing layup and manufacturing.}
\end{cvitems}
}
%------------------------------------------------
%------------------------------------------------
\end{cventries}